\begin{Info}
% Universidade
{UNIVERSIDADE DO VALE DO ITAJAÍ}
% Escola
%{ESCOLA DO MAR, CIÊNCIA E TECNOLOGIA}
{ESCOLA POLITÉCNICA}
% Curso
{CURSO DE CIÊNCIA DA COMPUTAÇÃO}
% Titulo
{Análise da Otimização de CNNs 1D na Classificação de Imagens Hiperespectrais}
% Autor
{Eduardo Savian}
% Cidade e Data
{Itajaí (SC), Dezembro de 2025}
% Nome da Área de concentração
{Visão Computacional}
% Orientador(a)
{Felipe Viel, PhD.}
% Coorientador(a) <Nome do Coorientador(a)>, <Titulação> %%%%%%%% Se não tiver coorientador deixe vazio
{}
\end{Info}

%\begin{Dedicatoria}
% Dedicatória
% \end{Dedicatoria}

% \begin{Agradecimentos}
% Agradeço a todos.
% \end{Agradecimentos}

\begin{Epigrafe}
Todos os Homens, por natureza, desejam conhecer - Aristóteles
\end{Epigrafe}

\begin{Resumo}
A classificação de imagens hiperespectrais (HSI) apresenta desafios relacionados à alta dimensionalidade dos dados e à complexidade da extração de informações espectrais discriminativas. Abordagens recentes que combinam redes neurais convolucionais unidimensionais (CNNs 1D) com técnicas de segmentação espacial têm demonstrado ganhos significativos de acurácia por meio da fusão espectral-espacial. No entanto, arquiteturas existentes apresentam limitações que incluem gargalos de processamento sequencial, ajuste manual de hiperparâmetros e exploração restrita de algoritmos de segmentação.

Este estudo experimental e quantitativo investiga o impacto de estratégias de otimização na classificação de HSI utilizando CNNs 1D combinadas com segmentação espacial. Experimentos sistemáticos são conduzidos em três conjuntos de dados de referência: Indian Pines, Salinas Valley e University of Pavia. A pesquisa avalia a paralelização das etapas de processamento, otimização automática de hiperparâmetros por meio de Grid Search e Otimização Bayesiana, e o desempenho comparativo dos algoritmos de segmentação SEEDS, LSC e Watershed.

As métricas de avaliação incluem Overall Accuracy, Average Accuracy, coeficiente Kappa, tempo de execução, consumo de memória e speedup. Os resultados esperados visam fornecer evidências empíricas sobre como essas estratégias de otimização afetam o desempenho de classificação e a eficiência computacional.

\textbf{Palavras-Chave}: Imagem Hiperespectral. Aprendizado de Máquina. Aprendizado Profundo. Segmentação. Rede Neural Convolucional Unidimensional. Otimização de Hiperparâmetros.
\end{Resumo}

\begin{Abstract}
\textit{Hyperspectral image (HSI) classification presents challenges related to high data dimensionality and the complexity of extracting discriminative spectral information. Recent approaches combining one-dimensional convolutional neural networks (1D CNNs) with spatial segmentation techniques have demonstrated significant accuracy gains through spectral-spatial fusion. However, existing architectures exhibit limitations including sequential processing bottlenecks, manual hyperparameter tuning, and restricted exploration of segmentation algorithms.}

\textit{This experimental and quantitative study investigates the impact of optimization strategies on HSI classification using 1D CNNs combined with spatial segmentation. Systematic experiments are conducted on three benchmark datasets: Indian Pines, Salinas Valley, and University of Pavia. The research evaluates parallelization of processing stages, automatic hyperparameter optimization through Grid Search and Bayesian Optimization, and comparative performance of SEEDS, LSC, and Watershed segmentation algorithms.}

\textit{Evaluation metrics include Overall Accuracy, Average Accuracy, Kappa coefficient, execution time, memory usage, and speedup. The expected findings aim to provide empirical evidence on how these optimization strategies affect classification performance and computational efficiency.}

\textit{\textbf{Keywords}: Hyperspectral Image. Machine Learning. Deep Learning. Segmentation. One-Dimensional Convolutional Neural Network. Hyperparameter Optimization.}
\end{Abstract}

